% Project 2 — Cache & Memory Performance Profiling
% Build: `latexmk` in this directory, or `make memory` from the repository root.
\documentclass[11pt]{article}
%% Packages and global settings shared by all ACS reports.
%% Loaded from each project's main.tex:  %% Packages and global settings shared by all ACS reports.
%% Loaded from each project's main.tex:  %% Packages and global settings shared by all ACS reports.
%% Loaded from each project's main.tex:  \input{../common/preamble}

\usepackage[T1]{fontenc}
\usepackage{lmodern}
\usepackage{microtype}
\usepackage[margin=1in]{geometry}
\usepackage{amsmath,amssymb}
\usepackage{graphicx}
\usepackage{booktabs,tabularx,multirow}
\usepackage{siunitx}
\usepackage{subcaption}
\usepackage[dvipsnames]{xcolor}
\usepackage{tcolorbox}
\tcbuselibrary{breakable}
\usepackage{listings}
\usepackage[shortlabels]{enumitem}
\usepackage{csquotes}

% Plots are drawn from the processed CSVs in <project>/data/ (synced from ACS-Code).
\usepackage{pgfplots}
\usepackage{pgfplotstable}
\pgfplotsset{compat=1.18}
\usepgfplotslibrary{groupplots}

\usepackage[backend=bibtex,style=ieee]{biblatex}
\addbibresource{references.bib}

\usepackage[hidelinks]{hyperref}
\usepackage[capitalise,noabbrev]{cleveref}

% ---------------------------------------------------------------- siunitx
\sisetup{
  per-mode = symbol,
  group-digits = integer,
  separate-uncertainty = true,
}
\DeclareSIUnit\flop{FLOP}
\DeclareSIUnit\iops{IOPS}
\DeclareSIUnit\op{op}
\DeclareSIUnit\elem{elem}
\DeclareSIUnit\cycle{cycle}
% Binary prefixes are built in: \qty{128}{\kibi\byte}, \qty{3.2}{\gibi\byte\per\second}

% ---------------------------------------------------------------- listings
\lstset{
  basicstyle = \ttfamily\footnotesize,
  columns = fullflexible,
  keepspaces = true,
  breaklines = true,
  frame = single,
  framerule = 0.3pt,
  rulecolor = \color{black!30},
  showstringspaces = false,
}
% Short disassembly / compiler-report excerpts:  \begin{lstlisting}[style=excerpt] ... \end{lstlisting}
\lstdefinestyle{excerpt}{basicstyle=\ttfamily\scriptsize, frame=leftline}


\usepackage[T1]{fontenc}
\usepackage{lmodern}
\usepackage{microtype}
\usepackage[margin=1in]{geometry}
\usepackage{amsmath,amssymb}
\usepackage{graphicx}
\usepackage{booktabs,tabularx,multirow}
\usepackage{siunitx}
\usepackage{subcaption}
\usepackage[dvipsnames]{xcolor}
\usepackage{tcolorbox}
\tcbuselibrary{breakable}
\usepackage{listings}
\usepackage[shortlabels]{enumitem}
\usepackage{csquotes}

% Plots are drawn from the processed CSVs in <project>/data/ (synced from ACS-Code).
\usepackage{pgfplots}
\usepackage{pgfplotstable}
\pgfplotsset{compat=1.18}
\usepgfplotslibrary{groupplots}

\usepackage[backend=bibtex,style=ieee]{biblatex}
\addbibresource{references.bib}

\usepackage[hidelinks]{hyperref}
\usepackage[capitalise,noabbrev]{cleveref}

% ---------------------------------------------------------------- siunitx
\sisetup{
  per-mode = symbol,
  group-digits = integer,
  separate-uncertainty = true,
}
\DeclareSIUnit\flop{FLOP}
\DeclareSIUnit\iops{IOPS}
\DeclareSIUnit\op{op}
\DeclareSIUnit\elem{elem}
\DeclareSIUnit\cycle{cycle}
% Binary prefixes are built in: \qty{128}{\kibi\byte}, \qty{3.2}{\gibi\byte\per\second}

% ---------------------------------------------------------------- listings
\lstset{
  basicstyle = \ttfamily\footnotesize,
  columns = fullflexible,
  keepspaces = true,
  breaklines = true,
  frame = single,
  framerule = 0.3pt,
  rulecolor = \color{black!30},
  showstringspaces = false,
}
% Short disassembly / compiler-report excerpts:  \begin{lstlisting}[style=excerpt] ... \end{lstlisting}
\lstdefinestyle{excerpt}{basicstyle=\ttfamily\scriptsize, frame=leftline}


\usepackage[T1]{fontenc}
\usepackage{lmodern}
\usepackage{microtype}
\usepackage[margin=1in]{geometry}
\usepackage{amsmath,amssymb}
\usepackage{graphicx}
\usepackage{booktabs,tabularx,multirow}
\usepackage{siunitx}
\usepackage{subcaption}
\usepackage[dvipsnames]{xcolor}
\usepackage{tcolorbox}
\tcbuselibrary{breakable}
\usepackage{listings}
\usepackage[shortlabels]{enumitem}
\usepackage{csquotes}

% Plots are drawn from the processed CSVs in <project>/data/ (synced from ACS-Code).
\usepackage{pgfplots}
\usepackage{pgfplotstable}
\pgfplotsset{compat=1.18}
\usepgfplotslibrary{groupplots}

\usepackage[backend=bibtex,style=ieee]{biblatex}
\addbibresource{references.bib}

\usepackage[hidelinks]{hyperref}
\usepackage[capitalise,noabbrev]{cleveref}

% ---------------------------------------------------------------- siunitx
\sisetup{
  per-mode = symbol,
  group-digits = integer,
  separate-uncertainty = true,
}
\DeclareSIUnit\flop{FLOP}
\DeclareSIUnit\iops{IOPS}
\DeclareSIUnit\op{op}
\DeclareSIUnit\elem{elem}
\DeclareSIUnit\cycle{cycle}
% Binary prefixes are built in: \qty{128}{\kibi\byte}, \qty{3.2}{\gibi\byte\per\second}

% ---------------------------------------------------------------- listings
\lstset{
  basicstyle = \ttfamily\footnotesize,
  columns = fullflexible,
  keepspaces = true,
  breaklines = true,
  frame = single,
  framerule = 0.3pt,
  rulecolor = \color{black!30},
  showstringspaces = false,
}
% Short disassembly / compiler-report excerpts:  \begin{lstlisting}[style=excerpt] ... \end{lstlisting}
\lstdefinestyle{excerpt}{basicstyle=\ttfamily\scriptsize, frame=leftline}

%% Macros shared by all ACS reports.  %% Macros shared by all ACS reports.  %% Macros shared by all ACS reports.  \input{../common/macros} after the preamble.

% ================================================================ draft / final
% Draft (default): rubric boxes are shown and missing data is drawn as a placeholder.
% Final (`make final`): rubric boxes are hidden; TODOs and missing data raise LaTeX warnings.
\newif\ifacsdraft
\acsdrafttrue
\ifdefined\acsfinal \acsdraftfalse \fi

% ================================================================ identity
\newcommand{\acsauthor}{Samson Kempiak}
\newcommand{\acsemail}{kempis@rpi.edu}
\newcommand{\acscourse}{ECSE 4320/6320 Advanced Computer Systems}
\newcommand{\acsdue}{October 21, 2026}
\newcommand{\coderepo}{\href{https://github.com/SKempiak/ACS-Code}{\texttt{github.com/SKempiak/ACS-Code}}}

% ================================================================ writing aids
% \todo{...}   something still to write
% \tbd         a number not measured yet
% \code{...}   file names, commands, flags (underscores are fine)
\newcommand{\todo}[1]{%
  \ifacsdraft\else\PackageWarning{acs}{Unresolved TODO}\fi
  \textcolor{red}{\textbf{[TODO:} #1\textbf{]}}}
\newcommand{\tbd}{%
  \ifacsdraft\else\PackageWarning{acs}{Unmeasured value (\string\tbd)}\fi
  \textcolor{orange}{\textbf{??}}}
\newcommand{\code}[1]{\texttt{\detokenize{#1}}}

% Rubric reminder at the top of a section (draft only):  \rubric{Category}{points}{what is graded}
\newtcolorbox{acsrubricbox}{colback=blue!4, colframe=blue!45, boxrule=0.4pt, arc=1pt,
  left=4pt, right=4pt, top=2pt, bottom=2pt, fontupper=\small}
\newcommand{\rubric}[3]{%
  \ifacsdraft\begin{acsrubricbox}\textbf{Rubric --- #1 (#2 pts).} #3\end{acsrubricbox}\fi}

% Pre-experiment hypothesis (common requirement): write and commit it before measuring.
\newtcolorbox{hypothesis}{breakable, colback=black!3, colframe=black!45, boxrule=0.4pt, arc=1pt,
  left=4pt, right=4pt, top=2pt, bottom=2pt, fonttitle=\bfseries\small,
  title={Hypothesis (from first principles, then compared with the data)}}

% Experiment setup table: one "Knob & value \\" row per line.
%   \begin{knobs}
%     Fixed  & kernel = SAXPY, float32 \\
%     Varied & $N = 2^{8} \ldots 2^{26}$ \\
%   \end{knobs}
\newenvironment{knobs}
  {\par\smallskip\noindent\small\tabularx{\linewidth}{@{}>{\bfseries}l X@{}}\toprule}
  {\bottomrule\endtabularx\par\smallskip}

% ================================================================ data & plots
% Processed CSVs live in <project>/data/<name>.csv (ACS-Code: scripts/sync_data.sh).
%
%   \IfDataExists{baseline}{<tikzpicture using data/baseline.csv>}
%
% draws the plot when the file exists, otherwise a placeholder box.
\newcommand{\acsmissing}[1]{%
  \ifacsdraft\else\PackageWarning{acs}{Missing data file #1}\fi
  \fbox{\parbox[c][0.22\textheight][c]{0.9\linewidth}{\centering\small\color{black!55}%
    Data pending: \texttt{\detokenize{#1}}\\[3pt]
    run the experiment in ACS-Code, then \texttt{scripts/sync\_data.sh}}}}
\newcommand{\IfDataExists}[2]{\IfFileExists{data/#1.csv}{#2}{\acsmissing{data/#1.csv}}}

\pgfplotsset{
  % Base style for every report plot.
  acs/.style={
    width=\linewidth, height=0.55\linewidth,
    grid=major, grid style={black!10},
    legend cell align=left, legend style={font=\footnotesize},
    label style={font=\small}, tick label style={font=\footnotesize},
    unbounded coords=discard,
    table/col sep=comma,
  },
  % Keep only rows where column #1 equals #2. Repeat the key to filter on several columns:
  %   \addplot table[x=n, y=ns_per_elem_median, only if={variant}{autovec}, only if={kernel}{saxpy}] {data/baseline.csv};
  /pgfplots/table/only if/.style 2 args={
    /pgfplots/x filter/.append code={%
      \edef\acsFilterValue{\thisrow{#1}}%
      \edef\acsFilterWant{#2}%
      \ifx\acsFilterValue\acsFilterWant\else\def\pgfmathresult{inf}\fi
    },
  },
  % Error bars spanning the interquartile range of a summarized value (tools/summarize columns):
  %   \addplot+[iqr bars] table[x=n, y=ns_per_elem_median, iqr of=ns_per_elem] {...};
  iqr bars/.style={error bars/.cd, y dir=both, y explicit},
  /pgfplots/table/iqr of/.style={
    y error plus expr={\thisrow{#1_q75}-\thisrow{#1_median}},
    y error minus expr={\thisrow{#1_median}-\thisrow{#1_q25}},
  },
}

% Vertical annotation inside an axis, e.g. cache-level boundaries:  \acsvline{32768}{L1d}
\newcommand{\acsvline}[2]{%
  \draw[densely dashed, black!50] ({axis cs:#1,1}|-{rel axis cs:0,0}) -- ({axis cs:#1,1}|-{rel axis cs:0,1})
    node[pos=0.98, anchor=east, rotate=90, font=\scriptsize, fill=white, inner sep=1pt] {#2};}
 after the preamble.

% ================================================================ draft / final
% Draft (default): rubric boxes are shown and missing data is drawn as a placeholder.
% Final (`make final`): rubric boxes are hidden; TODOs and missing data raise LaTeX warnings.
\newif\ifacsdraft
\acsdrafttrue
\ifdefined\acsfinal \acsdraftfalse \fi

% ================================================================ identity
\newcommand{\acsauthor}{Samson Kempiak}
\newcommand{\acsemail}{kempis@rpi.edu}
\newcommand{\acscourse}{ECSE 4320/6320 Advanced Computer Systems}
\newcommand{\acsdue}{October 21, 2026}
\newcommand{\coderepo}{\href{https://github.com/SKempiak/ACS-Code}{\texttt{github.com/SKempiak/ACS-Code}}}

% ================================================================ writing aids
% \todo{...}   something still to write
% \tbd         a number not measured yet
% \code{...}   file names, commands, flags (underscores are fine)
\newcommand{\todo}[1]{%
  \ifacsdraft\else\PackageWarning{acs}{Unresolved TODO}\fi
  \textcolor{red}{\textbf{[TODO:} #1\textbf{]}}}
\newcommand{\tbd}{%
  \ifacsdraft\else\PackageWarning{acs}{Unmeasured value (\string\tbd)}\fi
  \textcolor{orange}{\textbf{??}}}
\newcommand{\code}[1]{\texttt{\detokenize{#1}}}

% Rubric reminder at the top of a section (draft only):  \rubric{Category}{points}{what is graded}
\newtcolorbox{acsrubricbox}{colback=blue!4, colframe=blue!45, boxrule=0.4pt, arc=1pt,
  left=4pt, right=4pt, top=2pt, bottom=2pt, fontupper=\small}
\newcommand{\rubric}[3]{%
  \ifacsdraft\begin{acsrubricbox}\textbf{Rubric --- #1 (#2 pts).} #3\end{acsrubricbox}\fi}

% Pre-experiment hypothesis (common requirement): write and commit it before measuring.
\newtcolorbox{hypothesis}{breakable, colback=black!3, colframe=black!45, boxrule=0.4pt, arc=1pt,
  left=4pt, right=4pt, top=2pt, bottom=2pt, fonttitle=\bfseries\small,
  title={Hypothesis (from first principles, then compared with the data)}}

% Experiment setup table: one "Knob & value \\" row per line.
%   \begin{knobs}
%     Fixed  & kernel = SAXPY, float32 \\
%     Varied & $N = 2^{8} \ldots 2^{26}$ \\
%   \end{knobs}
\newenvironment{knobs}
  {\par\smallskip\noindent\small\tabularx{\linewidth}{@{}>{\bfseries}l X@{}}\toprule}
  {\bottomrule\endtabularx\par\smallskip}

% ================================================================ data & plots
% Processed CSVs live in <project>/data/<name>.csv (ACS-Code: scripts/sync_data.sh).
%
%   \IfDataExists{baseline}{<tikzpicture using data/baseline.csv>}
%
% draws the plot when the file exists, otherwise a placeholder box.
\newcommand{\acsmissing}[1]{%
  \ifacsdraft\else\PackageWarning{acs}{Missing data file #1}\fi
  \fbox{\parbox[c][0.22\textheight][c]{0.9\linewidth}{\centering\small\color{black!55}%
    Data pending: \texttt{\detokenize{#1}}\\[3pt]
    run the experiment in ACS-Code, then \texttt{scripts/sync\_data.sh}}}}
\newcommand{\IfDataExists}[2]{\IfFileExists{data/#1.csv}{#2}{\acsmissing{data/#1.csv}}}

\pgfplotsset{
  % Base style for every report plot.
  acs/.style={
    width=\linewidth, height=0.55\linewidth,
    grid=major, grid style={black!10},
    legend cell align=left, legend style={font=\footnotesize},
    label style={font=\small}, tick label style={font=\footnotesize},
    unbounded coords=discard,
    table/col sep=comma,
  },
  % Keep only rows where column #1 equals #2. Repeat the key to filter on several columns:
  %   \addplot table[x=n, y=ns_per_elem_median, only if={variant}{autovec}, only if={kernel}{saxpy}] {data/baseline.csv};
  /pgfplots/table/only if/.style 2 args={
    /pgfplots/x filter/.append code={%
      \edef\acsFilterValue{\thisrow{#1}}%
      \edef\acsFilterWant{#2}%
      \ifx\acsFilterValue\acsFilterWant\else\def\pgfmathresult{inf}\fi
    },
  },
  % Error bars spanning the interquartile range of a summarized value (tools/summarize columns):
  %   \addplot+[iqr bars] table[x=n, y=ns_per_elem_median, iqr of=ns_per_elem] {...};
  iqr bars/.style={error bars/.cd, y dir=both, y explicit},
  /pgfplots/table/iqr of/.style={
    y error plus expr={\thisrow{#1_q75}-\thisrow{#1_median}},
    y error minus expr={\thisrow{#1_median}-\thisrow{#1_q25}},
  },
}

% Vertical annotation inside an axis, e.g. cache-level boundaries:  \acsvline{32768}{L1d}
\newcommand{\acsvline}[2]{%
  \draw[densely dashed, black!50] ({axis cs:#1,1}|-{rel axis cs:0,0}) -- ({axis cs:#1,1}|-{rel axis cs:0,1})
    node[pos=0.98, anchor=east, rotate=90, font=\scriptsize, fill=white, inner sep=1pt] {#2};}
 after the preamble.

% ================================================================ draft / final
% Draft (default): rubric boxes are shown and missing data is drawn as a placeholder.
% Final (`make final`): rubric boxes are hidden; TODOs and missing data raise LaTeX warnings.
\newif\ifacsdraft
\acsdrafttrue
\ifdefined\acsfinal \acsdraftfalse \fi

% ================================================================ identity
\newcommand{\acsauthor}{Samson Kempiak}
\newcommand{\acsemail}{kempis@rpi.edu}
\newcommand{\acscourse}{ECSE 4320/6320 Advanced Computer Systems}
\newcommand{\acsdue}{October 21, 2026}
\newcommand{\coderepo}{\href{https://github.com/SKempiak/ACS-Code}{\texttt{github.com/SKempiak/ACS-Code}}}

% ================================================================ writing aids
% \todo{...}   something still to write
% \tbd         a number not measured yet
% \code{...}   file names, commands, flags (underscores are fine)
\newcommand{\todo}[1]{%
  \ifacsdraft\else\PackageWarning{acs}{Unresolved TODO}\fi
  \textcolor{red}{\textbf{[TODO:} #1\textbf{]}}}
\newcommand{\tbd}{%
  \ifacsdraft\else\PackageWarning{acs}{Unmeasured value (\string\tbd)}\fi
  \textcolor{orange}{\textbf{??}}}
\newcommand{\code}[1]{\texttt{\detokenize{#1}}}

% Rubric reminder at the top of a section (draft only):  \rubric{Category}{points}{what is graded}
\newtcolorbox{acsrubricbox}{colback=blue!4, colframe=blue!45, boxrule=0.4pt, arc=1pt,
  left=4pt, right=4pt, top=2pt, bottom=2pt, fontupper=\small}
\newcommand{\rubric}[3]{%
  \ifacsdraft\begin{acsrubricbox}\textbf{Rubric --- #1 (#2 pts).} #3\end{acsrubricbox}\fi}

% Pre-experiment hypothesis (common requirement): write and commit it before measuring.
\newtcolorbox{hypothesis}{breakable, colback=black!3, colframe=black!45, boxrule=0.4pt, arc=1pt,
  left=4pt, right=4pt, top=2pt, bottom=2pt, fonttitle=\bfseries\small,
  title={Hypothesis (from first principles, then compared with the data)}}

% Experiment setup table: one "Knob & value \\" row per line.
%   \begin{knobs}
%     Fixed  & kernel = SAXPY, float32 \\
%     Varied & $N = 2^{8} \ldots 2^{26}$ \\
%   \end{knobs}
\newenvironment{knobs}
  {\par\smallskip\noindent\small\tabularx{\linewidth}{@{}>{\bfseries}l X@{}}\toprule}
  {\bottomrule\endtabularx\par\smallskip}

% ================================================================ data & plots
% Processed CSVs live in <project>/data/<name>.csv (ACS-Code: scripts/sync_data.sh).
%
%   \IfDataExists{baseline}{<tikzpicture using data/baseline.csv>}
%
% draws the plot when the file exists, otherwise a placeholder box.
\newcommand{\acsmissing}[1]{%
  \ifacsdraft\else\PackageWarning{acs}{Missing data file #1}\fi
  \fbox{\parbox[c][0.22\textheight][c]{0.9\linewidth}{\centering\small\color{black!55}%
    Data pending: \texttt{\detokenize{#1}}\\[3pt]
    run the experiment in ACS-Code, then \texttt{scripts/sync\_data.sh}}}}
\newcommand{\IfDataExists}[2]{\IfFileExists{data/#1.csv}{#2}{\acsmissing{data/#1.csv}}}

\pgfplotsset{
  % Base style for every report plot.
  acs/.style={
    width=\linewidth, height=0.55\linewidth,
    grid=major, grid style={black!10},
    legend cell align=left, legend style={font=\footnotesize},
    label style={font=\small}, tick label style={font=\footnotesize},
    unbounded coords=discard,
    table/col sep=comma,
  },
  % Keep only rows where column #1 equals #2. Repeat the key to filter on several columns:
  %   \addplot table[x=n, y=ns_per_elem_median, only if={variant}{autovec}, only if={kernel}{saxpy}] {data/baseline.csv};
  /pgfplots/table/only if/.style 2 args={
    /pgfplots/x filter/.append code={%
      \edef\acsFilterValue{\thisrow{#1}}%
      \edef\acsFilterWant{#2}%
      \ifx\acsFilterValue\acsFilterWant\else\def\pgfmathresult{inf}\fi
    },
  },
  % Error bars spanning the interquartile range of a summarized value (tools/summarize columns):
  %   \addplot+[iqr bars] table[x=n, y=ns_per_elem_median, iqr of=ns_per_elem] {...};
  iqr bars/.style={error bars/.cd, y dir=both, y explicit},
  /pgfplots/table/iqr of/.style={
    y error plus expr={\thisrow{#1_q75}-\thisrow{#1_median}},
    y error minus expr={\thisrow{#1_median}-\thisrow{#1_q25}},
  },
}

% Vertical annotation inside an axis, e.g. cache-level boundaries:  \acsvline{32768}{L1d}
\newcommand{\acsvline}[2]{%
  \draw[densely dashed, black!50] ({axis cs:#1,1}|-{rel axis cs:0,0}) -- ({axis cs:#1,1}|-{rel axis cs:0,1})
    node[pos=0.98, anchor=east, rotate=90, font=\scriptsize, fill=white, inner sep=1pt] {#2};}

%% Measurement environment macros, used by common/environment_table.tex.
%% GENERATED by ACS-Code/scripts/collect_env.sh --tex on the test machine.
%% Regenerate with:
%%   scripts/collect_env.sh --tex > ../ACS-Latex/common/system_info.tex
\newcommand{\sysCPU}{Intel(R) Core(TM) i9-10900F CPU @ 2.80GHz}
\newcommand{\sysArch}{x86\_64}
\newcommand{\sysISA}{sse sse2 ssse3 sse4\_1 sse4\_2 avx avx2 fma f16c -- widest 256-bit (ymm)}
\newcommand{\sysCores}{1 socket(s), 10 cores/socket, 2 threads/core (20 logical)}
\newcommand{\sysSMT}{on (2 threads/core)}
\newcommand{\sysCaches}{L1d 320 KiB (10 instances); L2 2.5 MiB (10 instances); L3 20 MiB (1 instance)}
\newcommand{\sysNUMA}{1 node(s)}
\newcommand{\sysGovernor}{powersave}
\newcommand{\sysDriver}{intel\_pstate}
\newcommand{\sysTurbo}{enabled}
\newcommand{\sysEPP}{balance\_performance}
\newcommand{\sysIsolation}{no isolcpus; runs pinned with taskset/--cpu}
\newcommand{\sysMemory}{63 GiB total}
\newcommand{\sysDIMM}{not available (needs root/dmidecode)}
\newcommand{\sysPageSize}{4096 bytes}
\newcommand{\sysHugePages}{always}
\newcommand{\sysOS}{Arch Linux}
\newcommand{\sysKernel}{7.2.9-arch1-1}
\newcommand{\sysCC}{cc (GCC) 16.2.1 20260810}
\newcommand{\sysGCC}{gcc (GCC) 16.2.1 20260810}
\newcommand{\sysClang}{clang version 23.1.1}
\newcommand{\sysPerf}{not available (not installed)}
\newcommand{\sysPerfParanoid}{2}
\newcommand{\sysFio}{not available (not installed)}
\newcommand{\sysMLC}{not available (not installed)}
\newcommand{\sysGPU}{NVIDIA GeForce RTX 3070}
\newcommand{\sysCUDA}{not available (no CUDA)}
\newcommand{\sysCollected}{2026-10-12 18:35 EDT}


\title{Project 2: Cache \& Memory Performance Profiling}
\author{\acsauthor\\[2pt]\small\acscourse\\\small\href{mailto:\acsemail}{\acsemail}}
\date{\acsdue}

\begin{document}
\maketitle

\begin{abstract}
  We profile the cache and memory hierarchy of an Intel i9-10900F with three custom
  single-threaded/multi-threaded microbenchmarks. Zero-queue pointer-chase latency rises in clean steps
  from \ensuremath{\sim}\qty{1}{\nano\second} (L1, 4--5 cycles) through \ensuremath{\sim}\qty{4}{\nano\second} (L2) and the tens of
  cycles across the \qty{20}{\mebi\byte} L3 to \ensuremath{\sim}\qty{84}{\nano\second} (\ensuremath{\sim}\num{400} cycles) in
  DRAM. Single-core streaming bandwidth tops out at \ensuremath{\sim}\qty{14}{\gibi\byte\per\second} and scales to a
  controller-limited \ensuremath{\sim}\qty{30}{\gibi\byte\per\second} at a four-to-six-core knee that Little's Law
  attributes to \ensuremath{\sim}\num{42} outstanding cache lines. Random traversal costs up to \num{47}$\times$ more
  than sequential once the footprint leaves the LLC, quantifying the prefetcher's reach; huge pages
  showed no effect because the sandbox never granted them (a verified negative result). We close with the
  cache-reuse bandwidth pitfall---where a too-small buffer reports \qty{56}{\gibi\byte\per\second} of
  nonexistent ``DRAM'' bandwidth---and honest limitations from the absence of hardware counters.
\end{abstract}

\ifacsdraft\tableofcontents\clearpage\fi

% !TeX root = ../main.tex
\section{Introduction}\label{sec:intro}

\todo{Motivation: the L1/L2/L3/DRAM hierarchy on this machine (sizes, expected latencies) and why
locality, concurrency and page size decide real performance.}

This report answers:
\begin{enumerate}
  \item What is the zero-queue access latency of each level, and the DRAM bandwidth under different
        patterns, granularities and read/write mixes? (\cref{sec:zero-queue,sec:pattern,sec:rw-mix})
  \item Where does added concurrency stop buying throughput, and does Little's Law hold? (\cref{sec:intensity})
  \item How do working-set size, cache misses and TLB misses show up in runtime and counters?
        (\cref{sec:working-set,sec:cache-miss,sec:tlb})
  \item Which benchmark mistake produces misleading numbers? (\cref{sec:pitfall})
\end{enumerate}

\todo{One-paragraph summary of the main findings.}

% !TeX root = ../main.tex
\section{Methodology}\label{sec:method}
\rubric{Experimental rigor}{25}{Pinning and isolation, repetitions, documented settings and reproducible scripts.}

\subsection{Platform}
\newcommand{\acspinning}{single-thread runs and the loaded-latency probe pinned to CPU~2 (a physical
core, not CPU~0); multi-threaded runs place one thread per physical core}
\newcommand{\acsenvextra}{%
  Memory configuration & \sysMemory{} as reported by the kernel (\code{MemTotal}); single NUMA node;
  DIMM count and speed not probed (needs root/\code{dmidecode}) \\
}
All measurements were taken on the desktop system in \cref{tab:environment}: a ten-core Intel
i9-10900F (Comet Lake, Skylake-family core) with per-core \qty{32}{\kibi\byte} L1d and
\qty{256}{\kibi\byte} L2, a shared \qty{20}{\mebi\byte} L3, and DDR4 memory on one NUMA node. The table is
generated by \code{scripts/collect_env.sh} on the test host; each raw-data directory also stores its own
snapshot (\code{env.txt}) taken immediately before that experiment, so the settings of every figure are on
record even if the machine changes later.
%% Measurement environment table (common requirement: measurement discipline).
%% Values come from common/system_info.tex. Before \input-ing this file a project may define
%%   \acspinning   how runs were pinned/isolated
%%   \acsenvextra  extra rows, e.g.  \newcommand{\acsenvextra}{SSD & model, firmware, PCIe link \\}
\providecommand{\acspinning}{\todo{e.g.\ \code{taskset -c 2} on a P-core; isolcpus if used}}
\providecommand{\acsenvextra}{}
\begin{table}[tbp]
  \centering\small
  \caption{Measurement environment.}\label{tab:environment}
  \begin{tabularx}{\linewidth}{@{}>{\raggedright\arraybackslash}p{0.2\linewidth}X@{}}
    \toprule
    CPU               & \sysCPU{} (\sysArch) \\
    SIMD ISA          & \sysISA \\
    Topology          & \sysCores; NUMA: \sysNUMA \\
    SMT               & \sysSMT \\
    Caches            & \sysCaches \\
    Memory            & \sysMemory; DIMMs: \sysDIMM \\
    Pages             & base \sysPageSize; THP: \sysHugePages \\
    OS / kernel       & \sysOS; \sysKernel \\
    Frequency policy  & governor: \sysGovernor; driver: \sysDriver; turbo: \sysTurbo; EPP: \sysEPP \\
    Pinning / isolation & \acspinning; \sysIsolation \\
    Compilers         & \sysGCC; \sysClang \\
    Tools             & \sysPerf{} (perf\_event\_paranoid \sysPerfParanoid); \sysFio; MLC: \sysMLC \\
    GPU (optional)    & \sysGPU; CUDA \sysCUDA \\
    \acsenvextra
    Collected         & \sysCollected \\
    \bottomrule
  \end{tabularx}
\end{table}


\subsection{Tools}
All benchmarks share the harness in \code{ACS-Code/common} (monotonic raw timer, CPU pinning) plus a
project header \code{src/memprof.h} (page-policy allocation, ring generators, counters), and raw
per-repetition rows are reduced by \code{tools/summarize} (median and inter-quartile range).
\code{scripts/derive.py} then computes every derived number quoted in the text (plateaus, knee, Little's
Law, AMAT fit) into \code{summary.csv}; the report reads those values directly, so no number in the prose is
copied by hand.
\begin{knobs}
  \code{mem_latency}   & single-thread pointer-chase latency: working-set size, access size, node stride, page size, ring generator, timer, counters \\
  \code{mem_bandwidth} & streaming bandwidth with $T$ pinned threads: read/write/copy/triad/rwmix, access-size blocks, byte stride \\
  \code{mem_loaded}    & loaded latency: $T$ load-generator cores (stream, rwmix, or $k$ independent chases) plus a concurrent latency probe \\
  \code{kernel_miss}   & SAXPY loop whose footprint, order and page size vary, with counters \\
  Counters             & \code{perf_event_open(2)} called directly, user space only \cite{linux_perf} \\
  Intel MLC            & not installed; substitute cross-checks in \cref{sec:validation} \cite{intel_mlc} \\
\end{knobs}
The \code{perf} command-line tool is not installed and \code{perf_event_paranoid} is
\sysPerfParanoid{}, which still allows a process to count its \emph{own} user-mode events. The benchmarks
therefore open the counters themselves (\code{exclude_kernel=1}, \code{exclude_hv=1}): cycles,
instructions, L1D read misses, LLC read references and misses, and dTLB read misses (completed page walks).
Events are opened independently so the kernel may multiplex them; counts are scaled by
\code{time_enabled/time_running} and the smallest running fraction is recorded per row
(\code{ctr_cover}). Intel MLC is the spec's suggested cross-check but is not installed and could not be
installed for this study; \code{run.sh mlc} records that fact in its raw directory, and
\cref{sec:validation} describes the clearly labelled substitutes.

\subsection{Controls}
\begin{itemize}
  \item \textbf{Pinning and isolation:} \acspinning. SMT: \sysSMT. The machine is otherwise idle and
        experiments are run one at a time.
  \item \textbf{Frequency:} governor \sysGovernor, turbo \sysTurbo. Cycle figures use the effective
        core frequency each run measures in-process from a dependent add chain
        (\valqty{zq_ghz}{\giga\hertz} median), not a nominal clock.
  \item \textbf{Pages:} every buffer is a \qty{2}{\mebi\byte}-aligned anonymous mapping whose page size is
        chosen by \code{madvise} \emph{before first touch}: \code{MADV_HUGEPAGE} for ``huge'' and
        \code{MADV_NOHUGEPAGE} for ``4k''. The second is essential because THP is \sysHugePages{}: without
        it a ``\qty{4}{\kibi\byte}'' baseline can silently receive huge pages. Each run records the
        \code{AnonHugePages} actually granted (from \code{/proc/self/smaps_rollup}).
  \item \textbf{Pre-faulting and warm-up:} every buffer is written once to fault in its pages, then the
        ring or stream is traversed (or one untimed repetition is run) before timing, so neither page
        faults nor cold-cache effects fall inside the timed region.
  \item \textbf{Repetitions:} 8 per point; figures show the median with inter-quartile-range error bars.
  \item \textbf{Timing:} each timed interval is calibrated to $\geq$\qtyrange{40}{200}{\milli\second},
        far above the timer resolution, so timer overhead is negligible. The loaded-latency windows read
        generator progress counters at the start and end of the \emph{same} interval in which the probe
        runs, so throughput and latency describe the same moment.
  \item \textbf{Provenance:} each \code{data/raw/<experiment>/<timestamp>/} directory holds the raw rows,
        \code{env.txt}, \code{build.txt} (git commit, compiler, flags, THP and counter settings),
        \code{commands.txt} (every command line) and \code{validate.txt} (the sanity gate of
        \cref{sec:validation}, run immediately before the experiment).
\end{itemize}

% !TeX root = ../main.tex
\section{Validation and Sanity Checks}\label{sec:validation}

\todo{Checks that the tools measure what they claim, e.g.:}
\begin{itemize}
  \item \todo{L1 latency on the order of a few CPU cycles at the measured frequency.}
  \item \todo{Single-thread sequential bandwidth below the theoretical DIMM/channel limit.}
  \item \todo{Cross-check against Intel MLC or another reference for at least one configuration.}
  \item \todo{Timer and loop overhead measured with an empty access loop and subtracted or bounded.}
\end{itemize}

% !TeX root = ../main.tex
\section{Zero-Queue Baselines}\label{sec:zero-queue}
\rubric{Zero-queue baselines}{25}{Methodologically sound per-level latency measurements with clear units and caveats.}

\begin{hypothesis}
  \todo{Expected latency per level (in ns and in CPU cycles) and why.}
\end{hypothesis}

\begin{knobs}
  Fixed       & one thread pinned to CPU \tbd; random reads; granularity \tbd~bytes \\
  Varied      & working-set size placed in L1, L2, LLC, DRAM \\
  Script      & \code{Cache_and_Memory_Performance_Profiling/scripts/run.sh zero_queue} \\
\end{knobs}

% TODO: set to the measured operating frequency in GHz (used for the cycles column).
\newcommand{\cpuFreqGHz}{1.0}
\begin{table}[tbp]
  \centering\small
  \caption{Zero-queue access latency per level. Cycles assume \cpuFreqGHz~GHz \todo{measured frequency}.}\label{tab:zero-queue}
  \IfDataExists{zero_queue}{%
    \pgfplotstabletypeset[
      col sep=comma,
      columns={size_bytes,ns_per_access_median,ns_per_access_q25,ns_per_access_q75,cycles},
      create on use/cycles/.style={create col/expr={\thisrow{ns_per_access_median}*\cpuFreqGHz}},
      columns/size_bytes/.style={column name=Working set (bytes), int detect},
      columns/ns_per_access_median/.style={column name=Median (ns), fixed, precision=2},
      columns/ns_per_access_q25/.style={column name=Q1 (ns), fixed, precision=2},
      columns/ns_per_access_q75/.style={column name=Q3 (ns), fixed, precision=2},
      columns/cycles/.style={column name=Cycles, fixed, precision=1},
      every head row/.style={before row=\toprule, after row=\midrule},
      every last row/.style={after row=\bottomrule},
    ]{data/zero_queue.csv}%
  }
\end{table}

\subsection{Analysis}
\todo{Which level each size actually lands in (evidence), caveats (timer/loop overhead, TLB,
prefetchers), and comparison with published or MLC numbers.}

% !TeX root = ../main.tex
\section{Access Pattern and Granularity}\label{sec:pattern}
\rubric{Pattern, granularity and read/write sweeps}{35}{Complete, well-controlled bandwidth and latency
comparisons (with \cref{sec:rw-mix}).}

\begin{hypothesis}
  Sequential access should let the prefetcher amortise DRAM latency, giving high bandwidth that
  \emph{falls} as the stride grows and each fetched \qty{64}{\byte} line delivers fewer useful bytes.
  Random access should be near-flat and low across granularities, since every access is a fresh cache
  miss regardless of stride.
\end{hypothesis}

\begin{knobs}
  Fixed       & 1 thread pinned to CPU~2; \qty{256}{\mebi\byte} DRAM-resident region; 100\% reads \\
  Varied      & pattern (sequential, random); byte stride (8, 64, 256, 1024) \\
  Script      & \code{Cache_and_Memory_Performance_Profiling/scripts/run.sh granularity} and \code{latency} \\
\end{knobs}

\begin{figure}[tbp]
  \centering
  \begin{subfigure}[t]{0.49\linewidth}
    \IfDataExists{granularity}{%
    \begin{tikzpicture}
      \begin{axis}[acs, xmode=log, log basis x=2, ymode=log, height=0.85\linewidth,
                   xlabel={Byte stride}, ylabel={Bandwidth (GiB/s)}, legend pos=south west]
        \addplot+[iqr bars] table[x=strideb, y=gib_s_median, iqr of=gib_s, only if={pattern}{sequential}] {data/granularity.csv};
        \addlegendentry{sequential}
        \addplot+[iqr bars] table[x=strideb, y=gib_s_median, iqr of=gib_s, only if={pattern}{random}] {data/granularity.csv};
        \addlegendentry{random}
      \end{axis}
    \end{tikzpicture}}
    \caption{Delivered read bandwidth.}
  \end{subfigure}\hfill
  \begin{subfigure}[t]{0.49\linewidth}
    \IfDataExists{latency}{%
    \begin{tikzpicture}
      \begin{axis}[acs, xmode=log, log basis x=2, ymode=log, height=0.85\linewidth,
                   xlabel={Working set (bytes)}, ylabel={Latency (ns/access)}, legend pos=north west]
        \addplot+[iqr bars] table[x=bytes, y=ns_per_access_median, iqr of=ns_per_access, only if={pattern}{random}] {data/latency.csv};
        \addlegendentry{random}
        \addplot+[iqr bars] table[x=bytes, y=ns_per_access_median, iqr of=ns_per_access, only if={pattern}{sequential}] {data/latency.csv};
        \addlegendentry{sequential}
      \end{axis}
    \end{tikzpicture}}
    \caption{Pointer-chase latency vs.\ size.}
  \end{subfigure}%
  \caption{Access pattern and granularity. Left: sequential read bandwidth falls steeply with stride as
  each \qty{64}{\byte} line yields fewer useful bytes, while random bandwidth is flat and low. Right: the
  prefetcher keeps sequential latency near \qty{5}{\nano\second} even in DRAM, whereas random latency
  tracks the full hierarchy.}\label{fig:pattern}
\end{figure}

\subsection{Analysis}
Sequential bandwidth drops from \qty{14.8}{\gibi\byte\per\second} at unit (8-byte) stride to
\qty{2.4}{\gibi\byte\per\second} at \qty{64}{\byte}, \qty{1.09}{\gibi\byte\per\second} at
\qty{256}{\byte} and \qty{0.93}{\gibi\byte\per\second} at \qty{1024}{\byte}. The knee is at the
\qty{64}{\byte} cache line: once the stride reaches a line, every access still pulls a full line but uses
only one element of it, so \emph{useful} bandwidth falls roughly as $1/\text{stride}$ while the prefetcher
can no longer help. Random access is flat at \qtyrange{0.85}{0.93}{\gibi\byte\per\second} across all
strides---each access is an independent miss, so granularity barely matters. The latency panel is the
same story from the other side: sequential access to a DRAM-resident region costs \qty{5.3}{\nano\second}
(prefetched) while random access costs \qty{84}{\nano\second}, a \num{16}$\times$ gap that is entirely
the prefetcher. This is the mechanism behind the cache-miss results in \cref{sec:cache-miss}.

% !TeX root = ../main.tex
\section{Read/Write Mix}\label{sec:rw-mix}

\begin{hypothesis}
  \todo{Expected throughput and latency for 100\% read, 100\% write, 70/30 and 50/50, and why writes
  differ (read-for-ownership, cache pollution, write-combining).}
\end{hypothesis}

\begin{knobs}
  Fixed       & 1 thread; random access; granularity \tbd~bytes; region \tbd \\
  Varied      & read percentage: 100, 70, 50, 0 \\
  Script      & \code{Cache_and_Memory_Performance_Profiling/scripts/run.sh rw_mix} \\
\end{knobs}

\begin{figure}[tbp]
  \centering
  \IfDataExists{rw_mix}{%
    \begin{tikzpicture}
      \begin{axis}[acs, ybar, bar width=14pt, width=0.6\linewidth, xtick={0,50,70,100},
                   xlabel={Reads (\%)}, ylabel={Bandwidth (GiB/s)}, ymin=0, enlarge x limits=0.15]
        \addplot+[iqr bars] table[x=read_pct, y=gib_per_s_median, iqr of=gib_per_s] {data/rw_mix.csv};
      \end{axis}
    \end{tikzpicture}%
  }
  \caption{Bandwidth across read/write mixes. \todo{takeaway; add latency}}\label{fig:rw-mix}
\end{figure}

\subsection{Analysis}
\todo{Compare throughput and latency across mixes and explain the differences.}

% !TeX root = ../main.tex
\section{Intensity Sweep and the Throughput--Latency Knee}\label{sec:intensity}
\rubric{Throughput--latency knee}{35}{Clear knee identification, diminishing-returns discussion and a
Little's Law connection.}

\begin{hypothesis}
  \todo{Thread count at which throughput should stop scaling and latency should rise, and why
  (cores, memory channels, shared bandwidth).}
\end{hypothesis}

\begin{knobs}
  Fixed       & region \tbd~per thread; random reads; granularity \tbd~bytes \\
  Varied      & threads $T$ = 1, 2, 4, \ldots{} (pinned: \todo{CPU list}) \\
  Script      & \code{Cache_and_Memory_Performance_Profiling/scripts/run.sh intensity} \\
\end{knobs}

\begin{figure}[tbp]
  \centering
  \IfDataExists{intensity}{%
    \begin{tikzpicture}
      \begin{axis}[acs, width=0.75\linewidth, xlabel={Throughput (GiB/s)}, ylabel={Latency (ns/op)},
                   legend pos=north west, nodes near coords, point meta=explicit symbolic,
                   every node near coord/.append style={font=\scriptsize, anchor=south east}]
        \addplot+ table[x=gib_per_s_median, y=lat_mean_ns_median, meta=threads] {data/intensity.csv};
        \addlegendentry{mean}
        \addplot+ table[x=gib_per_s_median, y=lat_p99_ns_median, meta=threads] {data/intensity.csv};
        \addlegendentry{p99}
      \end{axis}
    \end{tikzpicture}%
  }
  \caption{Throughput vs.\ latency; labels give the thread count. \todo{mark the knee}}\label{fig:intensity}
\end{figure}

\subsection{Little's Law check}
With $T$ threads always busy, the number of requests in the system is $L = T$, so throughput $\lambda$
and mean latency $W$ should satisfy $L = \lambda W$ \cite{little1961proof}.

\begin{table}[tbp]
  \centering\small
  \caption{Little's Law: measured $\lambda W$ against the number of busy threads.}\label{tab:littles-law}
  \IfDataExists{intensity}{%
    \pgfplotstabletypeset[
      col sep=comma,
      columns={threads,ops_per_s_median,lat_mean_ns_median,lambdaW},
      create on use/lambdaW/.style={create col/expr={\thisrow{ops_per_s_median}*\thisrow{lat_mean_ns_median}/1e9}},
      columns/threads/.style={column name=$T$},
      columns/ops_per_s_median/.style={column name={$\lambda$ (ops/s)}, sci, precision=2},
      columns/lat_mean_ns_median/.style={column name={$W$ (ns)}, fixed, precision=2},
      columns/lambdaW/.style={column name={$L=\lambda W$}, fixed, precision=2},
      every head row/.style={before row=\toprule, after row=\midrule},
      every last row/.style={after row=\bottomrule},
    ]{data/intensity.csv}%
  }
\end{table}

\subsection{Analysis}
\todo{Where the knee is and why; diminishing returns beyond it; whether $\lambda W \approx T$ holds and
what a deviation means (sampling, think time, oversubscription).}

% !TeX root = ../main.tex
\section{Working-Set Sweep}\label{sec:working-set}
\rubric{Working-set, cache and TLB analysis}{35}{Counters and performance correlated correctly (with
\cref{sec:cache-miss,sec:tlb}).}

\begin{hypothesis}
  \todo{Where latency steps should appear given L1d = \tbd, L2 = \tbd, LLC = \tbd{} on this CPU.}
\end{hypothesis}

\begin{knobs}
  Fixed       & 1 thread pinned; random reads; granularity \tbd~bytes; base pages \\
  Varied      & working-set size from \qty{4}{\kibi\byte} to \tbd~\si{\gibi\byte} \\
  Script      & \code{Cache_and_Memory_Performance_Profiling/scripts/run.sh working_set} \\
\end{knobs}

\begin{figure}[tbp]
  \centering
  \IfDataExists{working_set}{%
    \begin{tikzpicture}
      \begin{axis}[acs, xmode=log, log basis x=2, ymode=log, xlabel={Working-set size (bytes)}, ylabel={Latency (ns/access)}]
        \addplot+[iqr bars] table[x=size_bytes, y=ns_per_access_median, iqr of=ns_per_access, only if={huge_requested}{0}] {data/working_set.csv};
        % TODO: cache boundaries in bytes, e.g. \acsvline{48000}{L1d} \acsvline{1280000}{L2} \acsvline{30000000}{LLC}
      \end{axis}
    \end{tikzpicture}%
  }
  \caption{Access latency across working-set sizes. \todo{takeaway}}\label{fig:working-set}
\end{figure}

\subsection{Analysis}
\todo{Transition regions, and how sharp or gradual they are and why (shared LLC, prefetchers, TLB reach).}

% !TeX root = ../main.tex
\section{Cache-Miss Impact}\label{sec:cache-miss}

\begin{hypothesis}
  \todo{How runtime per element should track cache misses as footprint and traversal order change.}
\end{hypothesis}

\begin{knobs}
  Kernel      & $z_i = x_i \cdot y_i$ (float64) \\
  Varied      & $N$ (footprint); traversal order (sequential, random permutation) \\
  Counters    & \todo{events used; note that \code{perf stat} covers the whole process} \\
  Script      & \code{PERF=1 Cache_and_Memory_Performance_Profiling/scripts/run.sh cache_miss} \\
\end{knobs}

\begin{figure}[tbp]
  \centering
  \IfDataExists{cache_miss}{%
    \begin{tikzpicture}
      \begin{axis}[acs, xmode=log, log basis x=2, ymode=log, xlabel={$N$ (elements)}, ylabel={Time per element (ns)}, legend pos=north west]
        \addplot+[iqr bars] table[x=n, y=ns_per_elem_median, iqr of=ns_per_elem, only if={order}{seq}] {data/cache_miss.csv};
        \addlegendentry{sequential}
        \addplot+[iqr bars] table[x=n, y=ns_per_elem_median, iqr of=ns_per_elem, only if={order}{random}] {data/cache_miss.csv};
        \addlegendentry{random}
        % TODO (PERF=1 data): misses per element on a second axis, e.g.
        % \addplot table[x=n, y expr={\thisrow{LLC-load-misses}/\thisrow{n}}, only if={order}{random}] {data/cache_miss.csv};
      \end{axis}
    \end{tikzpicture}%
  }
  \caption{Runtime per element versus footprint and order. \todo{takeaway}}\label{fig:cache-miss}
\end{figure}

\subsection{AMAT interpretation}
\begin{equation}
  \text{AMAT} = t_{L1} + m_{L1}\bigl(t_{L2} + m_{L2}\,(t_{L3} + m_{L3}\, t_{\text{DRAM}})\bigr)
\end{equation}
\todo{Plug in the zero-queue latencies (\cref{tab:zero-queue}) and measured local miss ratios; compare the
predicted and measured time per element.}

% !TeX root = ../main.tex
\section{TLB-Miss Impact}\label{sec:tlb}

\begin{hypothesis}
  \todo{Expected effect of 2~MiB pages on dTLB misses and runtime for random access over large
  footprints, given the TLB sizes of this CPU.}
\end{hypothesis}

\begin{knobs}
  Kernel      & $z_i = x_i \cdot y_i$ (float64), random traversal \\
  Varied      & page size (base, transparent huge pages); $N$ \\
  Verification & \todo{how huge-page backing was confirmed (AnonHugePages)} \\
  Script      & \code{PERF=1 Cache_and_Memory_Performance_Profiling/scripts/run.sh tlb} \\
\end{knobs}

\begin{figure}[tbp]
  \centering
  \IfDataExists{tlb}{%
    \begin{tikzpicture}
      \begin{axis}[acs, xmode=log, log basis x=2, xlabel={$N$ (elements)}, ylabel={Time per element (ns)}, legend pos=north west]
        \addplot+[iqr bars] table[x=n, y=ns_per_elem_median, iqr of=ns_per_elem, only if={huge_requested}{0}] {data/tlb.csv};
        \addlegendentry{base pages}
        \addplot+[iqr bars] table[x=n, y=ns_per_elem_median, iqr of=ns_per_elem, only if={huge_requested}{1}] {data/tlb.csv};
        \addlegendentry{huge pages}
      \end{axis}
    \end{tikzpicture}%
  }
  \caption{Runtime with base vs.\ huge pages. \todo{takeaway; add dTLB-load-misses}}\label{fig:tlb}
\end{figure}

\subsection{Analysis}
\todo{TLB reach (entries $\times$ page size) vs.\ footprint; change in dTLB misses and page walks;
runtime effect and its limits.}

% !TeX root = ../main.tex
\section{Benchmark Pitfall}\label{sec:pitfall}
\rubric{AI-use transparency and pitfall analysis}{15}{Honest AI-use note and a meaningful benchmark-pitfall
demonstration (with \cref{sec:ai-use}).}

\begin{hypothesis}
  \todo{What the flawed setup will report and why it looks plausible.}
\end{hypothesis}

\begin{knobs}
  Pitfall     & \todo{e.g.\ measuring cache reuse instead of DRAM, unpinned threads, frequency scaling} \\
  Wrong setup & \todo{} \\
  Right setup & \todo{} \\
  Script      & \code{Cache_and_Memory_Performance_Profiling/scripts/run.sh pitfall} \\
\end{knobs}

\todo{Figure or table comparing the two setups.}

\subsection{Explanation}
\todo{Why the flawed number is wrong, how to detect it, and how the methodology avoids it.}

% !TeX root = ../main.tex
\section{Anomalies and Limitations}\label{sec:anomalies}
\rubric{Report quality}{15}{Readable plots, concise interpretation and thoughtful limitations.}

\subsection{Unexpected results}
\begin{itemize}
  \item \todo{Observation $\rightarrow$ plausible hypothesis $\rightarrow$ how it was (or could be) tested.}
\end{itemize}

\subsection{Limitations}
\begin{itemize}
  \item \todo{e.g.\ single machine and memory configuration, sampled latency resolution, counters
        covering initialization, THP backing not guaranteed.}
\end{itemize}

% !TeX root = ../main.tex
\section{AI-Use Note}\label{sec:ai-use}

Summarized from \code{AI_USE_LOG.md} in the code repository.

\begin{description}[style=nextline]
  \item[Tools used] Claude Code (Claude Opus~5) for coding, debugging and drafting this report;
        GCC~\sysGCC{} for the builds. The AI wrote code and prose; every number comes from real runs on
        the machine in \cref{tab:environment}.
  \item[What they helped generate] The three microbenchmarks---\code{mem_latency} (Sattolo
        pointer-chase), \code{mem_bandwidth} (multi-thread streaming with read/write/copy/triad/rwmix
        modes) and \code{kernel_miss} (footprint-varying SAXPY)---the shared timer/allocation/pinning
        helpers, the experiment driver (\code{run.sh}) and the first drafts of these sections.
  \item[What was independently modified, tested, or rejected] Three concrete corrections.
        (1)~The first latency benchmark used a \emph{sequential} stride, which the prefetcher turned
        into a \ensuremath{\sim}\qty{5}{\nano\second} ``latency'' (really throughput); it was replaced with a random
        Sattolo cycle so exactly one request is outstanding---the sequential number is now kept only as
        a deliberate caveat row in \cref{tab:zero-queue}. (2)~The huge-page experiment initially
        \emph{assumed} THP backing because the governor was \code{always}; adding a runtime
        \code{/proc/self/smaps_rollup} check revealed \code{AnonHugePages: 0 kB}, turning a false
        ``huge pages don't help'' into an honest limitation (\cref{sec:tlb}). (3)~An AI-suggested
        bandwidth test defaulted to a small buffer and reported \qty{56}{\gibi\byte\per\second} as
        ``DRAM bandwidth''; this is the cache-reuse pitfall and is now documented as the required
        negative case (\cref{sec:pitfall}) rather than quoted as a result.
  \item[Validation] Sanity checks in \code{scripts/validate.sh} assert the physically necessary
        orderings (L1 $<$ L2 $<$ L3 $<$ DRAM latency; L1 bandwidth $>$ DRAM; multi-thread $>$
        single-thread) and all pass. The hierarchy shape is cross-checked against the known cache
        capacities; Intel MLC and \code{perf} counters were unavailable for an external cross-check and
        this is stated as a limitation (\cref{sec:anomalies}). Every figure traces to a committed
        raw/processed file (\cref{app:repro}).
\end{description}

% !TeX root = ../main.tex
\section{Conclusion}\label{sec:conclusion}

\todo{Three to five sentences: the latency and bandwidth picture of this hierarchy, where concurrency
stops paying off, and the most important lesson about locality, cache and TLB behavior.}


\printbibliography

\appendix
% !TeX root = ../main.tex
\section{Reproducibility}\label{app:repro}

Code, scripts, raw and processed data, and this report live in one repository with two top-level
directories, \code{ACS-Code/} and \code{ACS-Latex/}, at commit \code{2721691f277e}.

\begin{lstlisting}[language=bash]
gh repo clone SKempiak/ACS-Code ACS && cd ACS
git checkout 2721691f277e
cd ACS-Code
make tools                                            # tools/summarize
make -C Cache_and_Memory_Performance_Profiling        # mem_latency mem_bandwidth mem_loaded kernel_miss
make -C Cache_and_Memory_Performance_Profiling test   # physical-ordering sanity checks (validate.sh)
scripts/collect_env.sh --tex > ../ACS-Latex/common/system_info.tex
Cache_and_Memory_Performance_Profiling/scripts/run.sh zero_queue pattern rw_mix intensity \
    working_set cache_miss tlb pitfall mlc            # or: run.sh all (CPU_PIN=2 by default)
scripts/sync_data.sh Cache_and_Memory_Performance_Profiling
cd ../ACS-Latex && make memory
\end{lstlisting}

\begin{table}[h]
  \centering\small
  \caption{Raw data behind each figure and table (directories under
  \code{ACS-Code/Cache_and_Memory_Performance_Profiling/data/raw/}).}\label{tab:raw-data}
  \footnotesize
  \begin{tabularx}{\linewidth}{@{}>{\raggedright\arraybackslash}p{0.2\linewidth}l>{\raggedright\arraybackslash}Xl@{}}
    \toprule
    Figure / table & \code{run.sh} experiment & Processed CSV & Raw run directory \\
    \midrule
    \Cref{tab:zero-queue} & \code{zero_queue} & \code{zero_queue}, \code{zero_queue_table} & \code{zero_queue/20261021-004356} \\
    \Cref{fig:pattern,fig:stride} & \code{pattern} & \code{pattern_lat}, \code{pattern_bw} & \code{pattern/20261021-004356} \\
    \Cref{fig:rw-mix} & \code{rw_mix} & \code{rw_mix}, \code{rw_mix_ref} & \code{rw_mix/20261021-004356} \\
    \Cref{fig:intensity,fig:mlp}, \cref{tab:littles-law} & \code{intensity} & \code{intensity_loaded}, \code{intensity_mlp}, \code{littles_law}, \code{knee} & \code{intensity/20261021-004356} \\
    \Cref{fig:working-set} & \code{working_set} & \code{working_set} & \code{working_set/20261021-004356} \\
    \Cref{fig:cache-miss} & \code{cache_miss} & \code{cache_miss_ctr}, \code{amat_fit} & \code{cache_miss/20261021-004356} \\
    \Cref{fig:tlb} & \code{tlb} & \code{tlb_chase}, \code{tlb_saxpy} & \code{tlb/20261021-004356} \\
    \Cref{fig:pitfall} & \code{pitfall} & \code{pitfall} & \code{pitfall/20261021-004356} \\
    \Cref{sec:validation} (MLC) & \code{mlc} & --- & \code{mlc/20261021-004356} \\
    \bottomrule
  \end{tabularx}
\end{table}

Each raw directory holds the raw rows (one per repetition, with full configuration metadata),
\code{env.txt}, \code{build.txt}, \code{commands.txt} and \code{validate.txt} (\cref{sec:method}). The
processed CSVs are produced from them by \code{tools/summarize}; \code{scripts/derive.py} computes the
derived tables and \code{summary.csv}, from which every number in the text is typeset; and
\code{scripts/sync_data.sh} copies all of them into this report's \code{data/}. All nine directories come from a single \code{run.sh all} invocation
(2026-10-21, 00:43--00:55; one shared timestamp). The superseded runs of 2026-10-19 (\code{latency/},
\code{granularity/}, \code{rwmix/}, \code{intensity/}, \code{cache_miss/}, \code{pitfall/} and both
\code{tlb/} runs) and their processed CSVs were removed from the repository.


\end{document}
